% GENERATED FILE. Edit canonical lessons, then run npm run generate:books.
% canonical-source-hash: fnv1a64:717b06968341403d
% canonical-lessons: MR-C298-kincalne, MR-C298-gondhal, MR-C298-hahakar, MR-C298-jamav, MR-C298-kujbujne

\chapter{Uproar in the Street}
\label{ch:mr-uproar-in-the-street}
\hlchaptermodality{298}

\begin{chapteropening}
\textbf{By the end of this chapter:} \emph{I can describe a commotion: screams, a crowd, whispers.}

Complete the last lesson of chapter 298: I can describe a commotion: screams, a crowd, whispers.
\end{chapteropening}

\section[\texorpdfstring{\mr{किंचाळणे} (kiñcāḷṇe)}{किंचाळणे (kiñcāḷṇe)}]{\mr{किंचाळणे} (kiñcāḷṇe) — to scream; to shriek}

\label{lesson:MR-C298-kincalne}



\begin{quote}
Before the new one: say the Marathi for to loot, then the Marathi for to fleece.
\end{quote}

\subsection*{You'll want to know: \mr{किंचाळणे}}
\textbf{\mr{किंचाळणे}} — \emph{kiñcāḷṇe} — \textquotedblleft{}to scream; to shriek\textquotedblright{}.

\begin{quote}
\mr{साप} \mr{बघून} \mr{ती} \mr{जोरात} \mr{किंचाळली}.
\end{quote}

\begin{quote}
\emph{sāp baghūn tī jorāt kiñcāḷlī.}
\end{quote}

\begin{quote}
She screamed loudly when she saw the snake.
\end{quote}

Past (he …): \mr{तो} \mr{किंचाळला}, \emph{to kiñcāḷlā}.

Intransitive: \mr{जोरात} \mr{किंचाळणे}, to scream aloud.

\begin{cousinweb}
\textbf{\mr{किंचाळणे}} (\emph{kiñcāḷṇe}) is a Marathi verb of crying out; CDIAL gathers the Marathi word for a shriek under the cry-word \emph{*kīkkati}, 'screams' (CDIAL 3192), but the verb itself is not listed and its origin is not settled.
\end{cousinweb}

\begin{grammarlens}[title={What you've built}]
One more word for telling what happened.
\end{grammarlens}

\subsection*{Your turn}
\begin{itemize}
\raggedright
  \item \emph{Say it:} \emph{kiñcāḷṇe}
  \item \emph{Say it:} \emph{kiñcāḷṇe}, once more
  \item \emph{Say it:} \emph{sāp baghūn tī jorāt kiñcāḷlī.} — She screamed loudly when she saw the snake.
  \item \emph{Recall:} say the Marathi for to loot, then the Marathi for to fleece, then say \emph{kiñcāḷṇe} again
\end{itemize}

\begin{tcolorbox}[breakable,colback=teal!4,colframe=teal!35!black,title={Before you move on}]
What does \mr{किंचाळणे} mean? (\textquotedblleft{}To scream; to shriek\textquotedblright{}.) Say it once more.
\end{tcolorbox}
\section[\texorpdfstring{\mr{गोंधळ} (gondhaḷ)}{गोंधळ (gondhaḷ)}]{\mr{गोंधळ} (gondhaḷ) — commotion; confusion}

\label{lesson:MR-C298-gondhal}



\begin{quote}
Before the new one: say the Marathi for to fleece, then the Marathi for to scream.
\end{quote}

\subsection*{You'll want to know: \mr{गोंधळ}}
\textbf{\mr{गोंधळ}} — \emph{gondhaḷ} — \textquotedblleft{}commotion; confusion\textquotedblright{}.

\begin{quote}
\mr{वीज} \mr{गेल्यावर} \mr{सभागृहात} \mr{एकच} \mr{गोंधळ} \mr{झाला}.
\end{quote}

\begin{quote}
\emph{vīj gelyāvar sabhāgṛhāt ekac gondhaḷ jhālā.}
\end{quote}

\begin{quote}
When the power went out, the hall was in total commotion.
\end{quote}

Masculine.

\mr{गोंधळ} \mr{होणे}, a commotion breaks out; \mr{गोंधळात} \mr{पडणे}, to be confused.

\begin{cousinweb}
\textbf{\mr{गोंधळ}} (\emph{gondhaḷ}) is the Marathi word for an uproar or a muddle; its origin is not settled.
\end{cousinweb}

\begin{grammarlens}[title={What you've built}]
One more word for telling what happened.
\end{grammarlens}

\subsection*{Your turn}
\begin{itemize}
\raggedright
  \item \emph{Say it:} \emph{gondhaḷ}
  \item \emph{Say it:} \emph{gondhaḷ}, once more
  \item \emph{Say it:} \emph{vīj gelyāvar sabhāgṛhāt ekac gondhaḷ jhālā.} — When the power went out, the hall was in total commotion.
  \item \emph{Recall:} say the Marathi for to fleece, then the Marathi for to scream, then say \emph{gondhaḷ} again
\end{itemize}

\begin{tcolorbox}[breakable,colback=teal!4,colframe=teal!35!black,title={Before you move on}]
What does \mr{गोंधळ} mean? (\textquotedblleft{}Commotion; confusion\textquotedblright{}.) Say it once more.
\end{tcolorbox}
\section[\texorpdfstring{\mr{हाहाकार} (hāhākār)}{हाहाकार (hāhākār)}]{\mr{हाहाकार} (hāhākār) — havoc; an outcry}

\label{lesson:MR-C298-hahakar}



\begin{quote}
Before the new one: say the Marathi for to scream, then the Marathi for commotion.
\end{quote}

\subsection*{You'll want to know: \mr{हाहाकार}}
\textbf{\mr{हाहाकार}} — \emph{hāhākār} — \textquotedblleft{}havoc; an outcry\textquotedblright{}.

\begin{quote}
\mr{पुरामुळे} \mr{गावात} \mr{हाहाकार} \mr{माजला}.
\end{quote}

\begin{quote}
\emph{purāmuḷe gāvāt hāhākār mājlā.}
\end{quote}

\begin{quote}
The flood caused havoc in the village.
\end{quote}

Masculine.

\mr{हाहाकार} \mr{माजणे}, havoc breaks out.

\begin{cousinweb}
\textbf{\mr{हाहाकार}} (\emph{hāhākār}) is built from the cry of distress \emph{hā hā} and Sanskrit \emph{kāra}, 'a making, a doing' (Monier-Williams).
\end{cousinweb}

\begin{grammarlens}[title={What you've built}]
One more word for telling what happened.
\end{grammarlens}

\subsection*{Your turn}
\begin{itemize}
\raggedright
  \item \emph{Say it:} \emph{hāhākār}
  \item \emph{Say it:} \emph{hāhākār}, once more
  \item \emph{Say it:} \emph{purāmuḷe gāvāt hāhākār mājlā.} — The flood caused havoc in the village.
  \item \emph{Recall:} say the Marathi for to scream, then the Marathi for commotion, then say \emph{hāhākār} again
\end{itemize}

\begin{tcolorbox}[breakable,colback=teal!4,colframe=teal!35!black,title={Before you move on}]
What does \mr{हाहाकार} mean? (\textquotedblleft{}Havoc; an outcry\textquotedblright{}.) Say it once more.
\end{tcolorbox}
\section[\texorpdfstring{\mr{जमाव} (jamāv)}{जमाव (jamāv)}]{\mr{जमाव} (jamāv) — a mob; a gathered crowd}

\label{lesson:MR-C298-jamav}



\begin{quote}
Before the new one: say the Marathi for commotion, then the Marathi for havoc.
\end{quote}

\subsection*{You'll want to know: \mr{जमाव}}
\textbf{\mr{जमाव}} — \emph{jamāv} — \textquotedblleft{}a mob; a gathered crowd\textquotedblright{}.

\begin{quote}
\mr{अपघाताच्या} \mr{ठिकाणी} \mr{मोठा} \mr{जमाव} \mr{जमला}.
\end{quote}

\begin{quote}
\emph{apghātācyā ṭhikāṇī moṭhā jamāv jamlā.}
\end{quote}

\begin{quote}
A large crowd gathered at the scene of the accident.
\end{quote}

Masculine.

Built on \mr{जमणे}, to gather. \mr{गर्दी} is any crowd; \mr{जमाव} is one that has gathered round something.

\begin{cousinweb}
\textbf{\mr{जमाव}} (\emph{jamāv}) is built on \emph{jamṇe}, 'to gather', which CDIAL lists, as Marathi 'to collect', under Sanskrit \emph{yamyate} (CDIAL 10428).
\end{cousinweb}

\begin{grammarlens}[title={What you've built}]
One more word for telling what happened.
\end{grammarlens}

\subsection*{Your turn}
\begin{itemize}
\raggedright
  \item \emph{Say it:} \emph{jamāv}
  \item \emph{Say it:} \emph{jamāv}, once more
  \item \emph{Say it:} \emph{apghātācyā ṭhikāṇī moṭhā jamāv jamlā.} — A large crowd gathered at the scene of the accident.
  \item \emph{Recall:} say the Marathi for commotion, then the Marathi for havoc, then say \emph{jamāv} again
\end{itemize}

\begin{tcolorbox}[breakable,colback=teal!4,colframe=teal!35!black,title={Before you move on}]
What does \mr{जमाव} mean? (\textquotedblleft{}A mob; a gathered crowd\textquotedblright{}.) Say it once more.
\end{tcolorbox}
\section[\texorpdfstring{\mr{कुजबुजणे} (kujbujṇe)}{कुजबुजणे (kujbujṇe)}]{\mr{कुजबुजणे} (kujbujṇe) — to whisper}

\label{lesson:MR-C298-kujbujne}



\begin{quote}
Before the new one: say the Marathi for havoc, then the Marathi for a mob.
\end{quote}

\subsection*{You'll want to know: \mr{कुजबुजणे}}
\textbf{\mr{कुजबुजणे}} — \emph{kujbujṇe} — \textquotedblleft{}to whisper\textquotedblright{}.

\begin{quote}
\mr{वर्गात} \mr{दोन} \mr{मुली} \mr{हळू} \mr{आवाजात} \mr{कुजबुजत} \mr{होत्या}.
\end{quote}

\begin{quote}
\emph{vargāt don mulī haḷū āvājāt kujbujat hotyā.}
\end{quote}

\begin{quote}
Two girls were whispering softly in class.
\end{quote}

Past (he …): \mr{तो} \mr{कुजबुजला}, \emph{to kujbujlā}.

The listener takes -\mr{शी}: \mr{मित्राशी} \mr{कुजबुजणे}.

\begin{cousinweb}
\textbf{\mr{कुजबुजणे}} (\emph{kujbujṇe}) is a Marathi verb built on a doubled sound, \emph{kuj-buj}; its origin is not settled.
\end{cousinweb}

\begin{grammarlens}[title={What you've built}]
One more word for telling what happened.
\end{grammarlens}

\subsection*{Your turn}
\begin{itemize}
\raggedright
  \item \emph{Say it:} \emph{kujbujṇe}
  \item \emph{Say it:} \emph{kujbujṇe}, once more
  \item \emph{Say it:} \emph{vargāt don mulī haḷū āvājāt kujbujat hotyā.} — Two girls were whispering softly in class.
  \item \emph{Recall:} say the Marathi for havoc, then the Marathi for a mob, then say \emph{kujbujṇe} again
\end{itemize}

\begin{tcolorbox}[breakable,colback=teal!4,colframe=teal!35!black,title={Before you move on}]
What does \mr{कुजबुजणे} mean? (\textquotedblleft{}To whisper\textquotedblright{}.) Say it once more.
\end{tcolorbox}
